% paper91-copilot-research.tex — Copilot Research Advisor
%   Research assistance via Judgment Geometry: five advisor classes
%   that compose the ideation subsystem.
%   Paper 91 of the Judgment Geometry series.
% Compile: pdflatex paper91-copilot-research

\documentclass[11pt]{article}
\usepackage{lmodern}
% jugeo-common.tex — shared preamble for all Judgment Geometry papers
% Include via: % jugeo-common.tex — shared preamble for all Judgment Geometry papers
% Include via: % jugeo-common.tex — shared preamble for all Judgment Geometry papers
% Include via: \input{jugeo-common}

% ── Packages ────────────────────────────────────────────────────────────
\usepackage{amsmath,amssymb,amsthm,mathtools}
\usepackage{tikz-cd}
\usepackage{listings}
\usepackage{xcolor}
\usepackage{xspace}
\usepackage{booktabs}
\usepackage{multirow}
\usepackage{enumitem}
\usepackage[expansion=false]{microtype}
\usepackage{hyperref}
\usepackage{cleveref}
% \usepackage{thmtools}  % disabled: not available in this TeX installation
\usepackage{float}
\usepackage{caption}
\usepackage{subcaption}
\usepackage{tabularx}
\usepackage{stmaryrd}       % \llbracket, \rrbracket
\usepackage{bbm}            % \mathbbm{1}

% ── Page geometry ───────────────────────────────────────────────────────
\usepackage[margin=1in]{geometry}

% ── Theorem environments ────────────────────────────────────────────────
\theoremstyle{plain}
\newtheorem{theorem}{Theorem}[section]
\newtheorem{lemma}[theorem]{Lemma}
\newtheorem{proposition}[theorem]{Proposition}
\newtheorem{corollary}[theorem]{Corollary}

\theoremstyle{definition}
\newtheorem{definition}[theorem]{Definition}
\newtheorem{example}[theorem]{Example}
\newtheorem{construction}[theorem]{Construction}

\theoremstyle{remark}
\newtheorem{remark}[theorem]{Remark}
\newtheorem{notation}[theorem]{Notation}

% ── Python listings style ──────────────────────────────────────────────
\definecolor{jugeoBlue}{HTML}{2563EB}
\definecolor{jugeoGreen}{HTML}{059669}
\definecolor{jugeoRed}{HTML}{DC2626}
\definecolor{jugeoPurple}{HTML}{7C3AED}
\definecolor{jugeoGray}{HTML}{6B7280}
\definecolor{jugeoBg}{HTML}{F8FAFC}

\lstdefinestyle{jugeo-python}{
  language=Python,
  basicstyle=\ttfamily\small,
  keywordstyle=\color{jugeoBlue}\bfseries,
  stringstyle=\color{jugeoGreen},
  commentstyle=\color{jugeoGray}\itshape,
  numberstyle=\tiny\color{jugeoGray},
  backgroundcolor=\color{jugeoBg},
  numbers=left,
  numbersep=6pt,
  frame=single,
  framerule=0.4pt,
  rulecolor=\color{jugeoGray!40},
  breaklines=true,
  breakatwhitespace=true,
  showstringspaces=false,
  tabsize=4,
  xleftmargin=1.5em,
  framexleftmargin=1.5em,
  morekeywords={dataclass,frozen,field,Enum,IntEnum,auto,Any,
                Mapping,Sequence,Optional,match,case},
  literate={→}{{$\to$}}1 {←}{{$\leftarrow$}}1
           {≤}{{$\leq$}}1 {≥}{{$\geq$}}1 {≠}{{$\neq$}}1
           {∈}{{$\in$}}1 {∀}{{$\forall$}}1 {∃}{{$\exists$}}1
           {⊕}{{$\oplus$}}1 {⊖}{{$\ominus$}}1,
  escapeinside={(*@}{@*)},
}
\lstset{style=jugeo-python}

\lstdefinestyle{jugeo-output}{
  basicstyle=\ttfamily\small,
  backgroundcolor=\color{jugeoBg},
  frame=single,
  framerule=0.4pt,
  rulecolor=\color{jugeoGray!40},
  breaklines=true,
  showstringspaces=false,
  xleftmargin=1.5em,
  framexleftmargin=0.5em,
  numbers=none,
}

% ── JuGeo-specific macros ──────────────────────────────────────────────

% Judgment 8-tuple
\newcommand{\J}{\mathcal{J}}
\newcommand{\Jud}[8]{(#1,\,#2,\,#3,\,#4,\,#5,\,#6,\,#7,\,#8)}
\newcommand{\JudShort}[1]{\J_{#1}}

% Components
\newcommand{\coord}{\mathsf{c}}            % coordinate
\newcommand{\prop}{\varphi}                 % proposition
\newcommand{\carrier}{\mathcal{A}}          % carrier type
\newcommand{\evid}{\mathcal{E}}             % evidence bundle
\newcommand{\oblig}{\mathcal{O}}            % obligations
\newcommand{\obstr}{\mathcal{B}}            % obstructions
\newcommand{\trust}{\mathcal{T}}            % trust annotation
\newcommand{\prov}{\Pi}                     % provenance

% Trust algebra
\newcommand{\Talg}{\mathfrak{T}}            % trust ordered algebra
\newcommand{\Eadm}{\mathcal{E}_{\mathrm{adm}}}
\newcommand{\tleq}{\preceq}                 % trust ordering
\newcommand{\tjoin}{\oplus}                 % conservative join
\newcommand{\tatten}{\ominus}               % attenuation
\newcommand{\tpromote}{\uparrow_\pi}        % promotion
\newcommand{\tdemote}{\downarrow_\chi}       % demotion

% Trust tiers
\newcommand{\Tcontra}{\ensuremath{\mathsf{CONTRADICTED}}}
\newcommand{\Tunver}{\ensuremath{\mathsf{UNVERIFIED}}}
\newcommand{\Tcopilot}{\ensuremath{\mathsf{COPILOT}}}
\newcommand{\Toracle}{\ensuremath{\mathsf{ORACLE}}}
\newcommand{\Truntime}{\ensuremath{\mathsf{RUNTIME}}}
\newcommand{\Tsolver}{\ensuremath{\mathsf{SOLVER}}}
\newcommand{\Tproof}{\ensuremath{\mathsf{PROOF}}}

% Site and geometry
\newcommand{\Site}{\mathcal{S}}             % semantic site
\newcommand{\Cov}{\mathrm{Cov}}            % covering family
\newcommand{\Ob}{\mathrm{Ob}}              % objects
\newcommand{\Mor}{\mathrm{Mor}}            % morphisms
\newcommand{\res}{\mathrm{res}}            % restriction
\newcommand{\Cech}{\check{C}}              % Čech complex
\newcommand{\Hcoh}[1]{H^{#1}}             % cohomology group

% Descent
\newcommand{\descent}{\mathrm{desc}}
\newcommand{\glue}{\mathrm{glue}}
\newcommand{\overlap}[2]{#1 \cap #2}

% Semantic moves
\newcommand{\VERIFY}{\textsc{Verify}}
\newcommand{\CONSTRUCT}{\textsc{Construct}}
\newcommand{\NORMALIZE}{\textsc{Normalize}}
\newcommand{\GLUE}{\textsc{Glue}}
\newcommand{\RESTRICT}{\textsc{Restrict}}
\newcommand{\TRANSPORT}{\textsc{Transport}}
\newcommand{\REPAIR}{\textsc{Repair}}
\newcommand{\PROMOTE}{\textsc{Promote}}
\newcommand{\CHALLENGE}{\textsc{Challenge}}
\newcommand{\ESCALATE}{\textsc{Escalate}}

% Fragments
\newcommand{\QFLIA}{\texttt{QF\_LIA}}
\newcommand{\QFLRA}{\texttt{QF\_LRA}}
\newcommand{\QFBV}{\texttt{QF\_BV}}
\newcommand{\QFUF}{\texttt{QF\_UF}}

% Misc
\newcommand{\Python}{\textsc{Python}}
\newcommand{\JuGeo}{\textsc{JuGeo}}
\newcommand{\LEAN}{\textsc{Lean}\,4}
\newcommand{\Fstar}{F$^{\star}$}
\newcommand{\Coq}{\textsc{Coq}}
\newcommand{\Dafny}{\textsc{Dafny}}

% Common shortcuts
\newcommand{\ie}{i.e.\@\xspace}
\newcommand{\eg}{e.g.\@\xspace}
\newcommand{\cf}{cf.\@\xspace}
\newcommand{\wrt}{w.r.t.\@\xspace}

% Evaluation shortcuts
\newcommand{\numcases}{300}
\newcommand{\numspec}{100}
\newcommand{\numeq}{100}
\newcommand{\numbug}{100}
\newcommand{\medtime}{6.5\,ms}
\newcommand{\totaltime}{1.949\,s}


% ── Packages ────────────────────────────────────────────────────────────
\usepackage{amsmath,amssymb,amsthm,mathtools}
\usepackage{tikz-cd}
\usepackage{listings}
\usepackage{xcolor}
\usepackage{xspace}
\usepackage{booktabs}
\usepackage{multirow}
\usepackage{enumitem}
\usepackage[expansion=false]{microtype}
\usepackage{hyperref}
\usepackage{cleveref}
% \usepackage{thmtools}  % disabled: not available in this TeX installation
\usepackage{float}
\usepackage{caption}
\usepackage{subcaption}
\usepackage{tabularx}
\usepackage{stmaryrd}       % \llbracket, \rrbracket
\usepackage{bbm}            % \mathbbm{1}

% ── Page geometry ───────────────────────────────────────────────────────
\usepackage[margin=1in]{geometry}

% ── Theorem environments ────────────────────────────────────────────────
\theoremstyle{plain}
\newtheorem{theorem}{Theorem}[section]
\newtheorem{lemma}[theorem]{Lemma}
\newtheorem{proposition}[theorem]{Proposition}
\newtheorem{corollary}[theorem]{Corollary}

\theoremstyle{definition}
\newtheorem{definition}[theorem]{Definition}
\newtheorem{example}[theorem]{Example}
\newtheorem{construction}[theorem]{Construction}

\theoremstyle{remark}
\newtheorem{remark}[theorem]{Remark}
\newtheorem{notation}[theorem]{Notation}

% ── Python listings style ──────────────────────────────────────────────
\definecolor{jugeoBlue}{HTML}{2563EB}
\definecolor{jugeoGreen}{HTML}{059669}
\definecolor{jugeoRed}{HTML}{DC2626}
\definecolor{jugeoPurple}{HTML}{7C3AED}
\definecolor{jugeoGray}{HTML}{6B7280}
\definecolor{jugeoBg}{HTML}{F8FAFC}

\lstdefinestyle{jugeo-python}{
  language=Python,
  basicstyle=\ttfamily\small,
  keywordstyle=\color{jugeoBlue}\bfseries,
  stringstyle=\color{jugeoGreen},
  commentstyle=\color{jugeoGray}\itshape,
  numberstyle=\tiny\color{jugeoGray},
  backgroundcolor=\color{jugeoBg},
  numbers=left,
  numbersep=6pt,
  frame=single,
  framerule=0.4pt,
  rulecolor=\color{jugeoGray!40},
  breaklines=true,
  breakatwhitespace=true,
  showstringspaces=false,
  tabsize=4,
  xleftmargin=1.5em,
  framexleftmargin=1.5em,
  morekeywords={dataclass,frozen,field,Enum,IntEnum,auto,Any,
                Mapping,Sequence,Optional,match,case},
  literate={→}{{$\to$}}1 {←}{{$\leftarrow$}}1
           {≤}{{$\leq$}}1 {≥}{{$\geq$}}1 {≠}{{$\neq$}}1
           {∈}{{$\in$}}1 {∀}{{$\forall$}}1 {∃}{{$\exists$}}1
           {⊕}{{$\oplus$}}1 {⊖}{{$\ominus$}}1,
  escapeinside={(*@}{@*)},
}
\lstset{style=jugeo-python}

\lstdefinestyle{jugeo-output}{
  basicstyle=\ttfamily\small,
  backgroundcolor=\color{jugeoBg},
  frame=single,
  framerule=0.4pt,
  rulecolor=\color{jugeoGray!40},
  breaklines=true,
  showstringspaces=false,
  xleftmargin=1.5em,
  framexleftmargin=0.5em,
  numbers=none,
}

% ── JuGeo-specific macros ──────────────────────────────────────────────

% Judgment 8-tuple
\newcommand{\J}{\mathcal{J}}
\newcommand{\Jud}[8]{(#1,\,#2,\,#3,\,#4,\,#5,\,#6,\,#7,\,#8)}
\newcommand{\JudShort}[1]{\J_{#1}}

% Components
\newcommand{\coord}{\mathsf{c}}            % coordinate
\newcommand{\prop}{\varphi}                 % proposition
\newcommand{\carrier}{\mathcal{A}}          % carrier type
\newcommand{\evid}{\mathcal{E}}             % evidence bundle
\newcommand{\oblig}{\mathcal{O}}            % obligations
\newcommand{\obstr}{\mathcal{B}}            % obstructions
\newcommand{\trust}{\mathcal{T}}            % trust annotation
\newcommand{\prov}{\Pi}                     % provenance

% Trust algebra
\newcommand{\Talg}{\mathfrak{T}}            % trust ordered algebra
\newcommand{\Eadm}{\mathcal{E}_{\mathrm{adm}}}
\newcommand{\tleq}{\preceq}                 % trust ordering
\newcommand{\tjoin}{\oplus}                 % conservative join
\newcommand{\tatten}{\ominus}               % attenuation
\newcommand{\tpromote}{\uparrow_\pi}        % promotion
\newcommand{\tdemote}{\downarrow_\chi}       % demotion

% Trust tiers
\newcommand{\Tcontra}{\ensuremath{\mathsf{CONTRADICTED}}}
\newcommand{\Tunver}{\ensuremath{\mathsf{UNVERIFIED}}}
\newcommand{\Tcopilot}{\ensuremath{\mathsf{COPILOT}}}
\newcommand{\Toracle}{\ensuremath{\mathsf{ORACLE}}}
\newcommand{\Truntime}{\ensuremath{\mathsf{RUNTIME}}}
\newcommand{\Tsolver}{\ensuremath{\mathsf{SOLVER}}}
\newcommand{\Tproof}{\ensuremath{\mathsf{PROOF}}}

% Site and geometry
\newcommand{\Site}{\mathcal{S}}             % semantic site
\newcommand{\Cov}{\mathrm{Cov}}            % covering family
\newcommand{\Ob}{\mathrm{Ob}}              % objects
\newcommand{\Mor}{\mathrm{Mor}}            % morphisms
\newcommand{\res}{\mathrm{res}}            % restriction
\newcommand{\Cech}{\check{C}}              % Čech complex
\newcommand{\Hcoh}[1]{H^{#1}}             % cohomology group

% Descent
\newcommand{\descent}{\mathrm{desc}}
\newcommand{\glue}{\mathrm{glue}}
\newcommand{\overlap}[2]{#1 \cap #2}

% Semantic moves
\newcommand{\VERIFY}{\textsc{Verify}}
\newcommand{\CONSTRUCT}{\textsc{Construct}}
\newcommand{\NORMALIZE}{\textsc{Normalize}}
\newcommand{\GLUE}{\textsc{Glue}}
\newcommand{\RESTRICT}{\textsc{Restrict}}
\newcommand{\TRANSPORT}{\textsc{Transport}}
\newcommand{\REPAIR}{\textsc{Repair}}
\newcommand{\PROMOTE}{\textsc{Promote}}
\newcommand{\CHALLENGE}{\textsc{Challenge}}
\newcommand{\ESCALATE}{\textsc{Escalate}}

% Fragments
\newcommand{\QFLIA}{\texttt{QF\_LIA}}
\newcommand{\QFLRA}{\texttt{QF\_LRA}}
\newcommand{\QFBV}{\texttt{QF\_BV}}
\newcommand{\QFUF}{\texttt{QF\_UF}}

% Misc
\newcommand{\Python}{\textsc{Python}}
\newcommand{\JuGeo}{\textsc{JuGeo}}
\newcommand{\LEAN}{\textsc{Lean}\,4}
\newcommand{\Fstar}{F$^{\star}$}
\newcommand{\Coq}{\textsc{Coq}}
\newcommand{\Dafny}{\textsc{Dafny}}

% Common shortcuts
\newcommand{\ie}{i.e.\@\xspace}
\newcommand{\eg}{e.g.\@\xspace}
\newcommand{\cf}{cf.\@\xspace}
\newcommand{\wrt}{w.r.t.\@\xspace}

% Evaluation shortcuts
\newcommand{\numcases}{300}
\newcommand{\numspec}{100}
\newcommand{\numeq}{100}
\newcommand{\numbug}{100}
\newcommand{\medtime}{6.5\,ms}
\newcommand{\totaltime}{1.949\,s}


% ── Packages ────────────────────────────────────────────────────────────
\usepackage{amsmath,amssymb,amsthm,mathtools}
\usepackage{tikz-cd}
\usepackage{listings}
\usepackage{xcolor}
\usepackage{xspace}
\usepackage{booktabs}
\usepackage{multirow}
\usepackage{enumitem}
\usepackage[expansion=false]{microtype}
\usepackage{hyperref}
\usepackage{cleveref}
% \usepackage{thmtools}  % disabled: not available in this TeX installation
\usepackage{float}
\usepackage{caption}
\usepackage{subcaption}
\usepackage{tabularx}
\usepackage{stmaryrd}       % \llbracket, \rrbracket
\usepackage{bbm}            % \mathbbm{1}

% ── Page geometry ───────────────────────────────────────────────────────
\usepackage[margin=1in]{geometry}

% ── Theorem environments ────────────────────────────────────────────────
\theoremstyle{plain}
\newtheorem{theorem}{Theorem}[section]
\newtheorem{lemma}[theorem]{Lemma}
\newtheorem{proposition}[theorem]{Proposition}
\newtheorem{corollary}[theorem]{Corollary}

\theoremstyle{definition}
\newtheorem{definition}[theorem]{Definition}
\newtheorem{example}[theorem]{Example}
\newtheorem{construction}[theorem]{Construction}

\theoremstyle{remark}
\newtheorem{remark}[theorem]{Remark}
\newtheorem{notation}[theorem]{Notation}

% ── Python listings style ──────────────────────────────────────────────
\definecolor{jugeoBlue}{HTML}{2563EB}
\definecolor{jugeoGreen}{HTML}{059669}
\definecolor{jugeoRed}{HTML}{DC2626}
\definecolor{jugeoPurple}{HTML}{7C3AED}
\definecolor{jugeoGray}{HTML}{6B7280}
\definecolor{jugeoBg}{HTML}{F8FAFC}

\lstdefinestyle{jugeo-python}{
  language=Python,
  basicstyle=\ttfamily\small,
  keywordstyle=\color{jugeoBlue}\bfseries,
  stringstyle=\color{jugeoGreen},
  commentstyle=\color{jugeoGray}\itshape,
  numberstyle=\tiny\color{jugeoGray},
  backgroundcolor=\color{jugeoBg},
  numbers=left,
  numbersep=6pt,
  frame=single,
  framerule=0.4pt,
  rulecolor=\color{jugeoGray!40},
  breaklines=true,
  breakatwhitespace=true,
  showstringspaces=false,
  tabsize=4,
  xleftmargin=1.5em,
  framexleftmargin=1.5em,
  morekeywords={dataclass,frozen,field,Enum,IntEnum,auto,Any,
                Mapping,Sequence,Optional,match,case},
  literate={→}{{$\to$}}1 {←}{{$\leftarrow$}}1
           {≤}{{$\leq$}}1 {≥}{{$\geq$}}1 {≠}{{$\neq$}}1
           {∈}{{$\in$}}1 {∀}{{$\forall$}}1 {∃}{{$\exists$}}1
           {⊕}{{$\oplus$}}1 {⊖}{{$\ominus$}}1,
  escapeinside={(*@}{@*)},
}
\lstset{style=jugeo-python}

\lstdefinestyle{jugeo-output}{
  basicstyle=\ttfamily\small,
  backgroundcolor=\color{jugeoBg},
  frame=single,
  framerule=0.4pt,
  rulecolor=\color{jugeoGray!40},
  breaklines=true,
  showstringspaces=false,
  xleftmargin=1.5em,
  framexleftmargin=0.5em,
  numbers=none,
}

% ── JuGeo-specific macros ──────────────────────────────────────────────

% Judgment 8-tuple
\newcommand{\J}{\mathcal{J}}
\newcommand{\Jud}[8]{(#1,\,#2,\,#3,\,#4,\,#5,\,#6,\,#7,\,#8)}
\newcommand{\JudShort}[1]{\J_{#1}}

% Components
\newcommand{\coord}{\mathsf{c}}            % coordinate
\newcommand{\prop}{\varphi}                 % proposition
\newcommand{\carrier}{\mathcal{A}}          % carrier type
\newcommand{\evid}{\mathcal{E}}             % evidence bundle
\newcommand{\oblig}{\mathcal{O}}            % obligations
\newcommand{\obstr}{\mathcal{B}}            % obstructions
\newcommand{\trust}{\mathcal{T}}            % trust annotation
\newcommand{\prov}{\Pi}                     % provenance

% Trust algebra
\newcommand{\Talg}{\mathfrak{T}}            % trust ordered algebra
\newcommand{\Eadm}{\mathcal{E}_{\mathrm{adm}}}
\newcommand{\tleq}{\preceq}                 % trust ordering
\newcommand{\tjoin}{\oplus}                 % conservative join
\newcommand{\tatten}{\ominus}               % attenuation
\newcommand{\tpromote}{\uparrow_\pi}        % promotion
\newcommand{\tdemote}{\downarrow_\chi}       % demotion

% Trust tiers
\newcommand{\Tcontra}{\ensuremath{\mathsf{CONTRADICTED}}}
\newcommand{\Tunver}{\ensuremath{\mathsf{UNVERIFIED}}}
\newcommand{\Tcopilot}{\ensuremath{\mathsf{COPILOT}}}
\newcommand{\Toracle}{\ensuremath{\mathsf{ORACLE}}}
\newcommand{\Truntime}{\ensuremath{\mathsf{RUNTIME}}}
\newcommand{\Tsolver}{\ensuremath{\mathsf{SOLVER}}}
\newcommand{\Tproof}{\ensuremath{\mathsf{PROOF}}}

% Site and geometry
\newcommand{\Site}{\mathcal{S}}             % semantic site
\newcommand{\Cov}{\mathrm{Cov}}            % covering family
\newcommand{\Ob}{\mathrm{Ob}}              % objects
\newcommand{\Mor}{\mathrm{Mor}}            % morphisms
\newcommand{\res}{\mathrm{res}}            % restriction
\newcommand{\Cech}{\check{C}}              % Čech complex
\newcommand{\Hcoh}[1]{H^{#1}}             % cohomology group

% Descent
\newcommand{\descent}{\mathrm{desc}}
\newcommand{\glue}{\mathrm{glue}}
\newcommand{\overlap}[2]{#1 \cap #2}

% Semantic moves
\newcommand{\VERIFY}{\textsc{Verify}}
\newcommand{\CONSTRUCT}{\textsc{Construct}}
\newcommand{\NORMALIZE}{\textsc{Normalize}}
\newcommand{\GLUE}{\textsc{Glue}}
\newcommand{\RESTRICT}{\textsc{Restrict}}
\newcommand{\TRANSPORT}{\textsc{Transport}}
\newcommand{\REPAIR}{\textsc{Repair}}
\newcommand{\PROMOTE}{\textsc{Promote}}
\newcommand{\CHALLENGE}{\textsc{Challenge}}
\newcommand{\ESCALATE}{\textsc{Escalate}}

% Fragments
\newcommand{\QFLIA}{\texttt{QF\_LIA}}
\newcommand{\QFLRA}{\texttt{QF\_LRA}}
\newcommand{\QFBV}{\texttt{QF\_BV}}
\newcommand{\QFUF}{\texttt{QF\_UF}}

% Misc
\newcommand{\Python}{\textsc{Python}}
\newcommand{\JuGeo}{\textsc{JuGeo}}
\newcommand{\LEAN}{\textsc{Lean}\,4}
\newcommand{\Fstar}{F$^{\star}$}
\newcommand{\Coq}{\textsc{Coq}}
\newcommand{\Dafny}{\textsc{Dafny}}

% Common shortcuts
\newcommand{\ie}{i.e.\@\xspace}
\newcommand{\eg}{e.g.\@\xspace}
\newcommand{\cf}{cf.\@\xspace}
\newcommand{\wrt}{w.r.t.\@\xspace}

% Evaluation shortcuts
\newcommand{\numcases}{300}
\newcommand{\numspec}{100}
\newcommand{\numeq}{100}
\newcommand{\numbug}{100}
\newcommand{\medtime}{6.5\,ms}
\newcommand{\totaltime}{1.949\,s}

% Auto-generated by experiments/exp91_copilot_research.py
% DO NOT EDIT — regenerate with: python3 experiments/exp91_copilot_research.py
%
% Macro prefix: \ppXCI   (Paper 91 — Copilot Research Advisor)

% ── Research Advisor metrics ────────────────────────────────────
\newcommand{\ppXCIsuggestionsGenerated}{383}
\newcommand{\ppXCIverificationPassRate}{78.9\%}
\newcommand{\ppXCImeanSuggestionLatency}{0.00\,s}
\newcommand{\ppXCIsoundSuggestions}{302}
\newcommand{\ppXCIoracleQueries}{1\,201}
\newcommand{\ppXCIoracleHitRate}{75.9\%}

% ── Experiment Advisor metrics ──────────────────────────────────
\newcommand{\ppXCIdesignsEvaluated}{128}
\newcommand{\ppXCIdesignImprovements}{95}
\newcommand{\ppXCIdesignImprovementRate}{74.2\%}
\newcommand{\ppXCImeanPowerGain}{0.20}
\newcommand{\ppXCIinsightsGenerated}{248}
\newcommand{\ppXCIpriorityReorderings}{12}

% ── Optimization Advisor metrics ────────────────────────────────
\newcommand{\ppXCIoptIterations}{512}
\newcommand{\ppXCIoptMonotonicRuns}{457}
\newcommand{\ppXCIoptMonotonicity}{89.3\%}
\newcommand{\ppXCIoptParetoPoints}{1\,599}
\newcommand{\ppXCImeanFrontImprovement}{1.2\%}
\newcommand{\ppXCIoptNextAccepted}{82.2\%}

% ── Futures Advisor metrics ─────────────────────────────────────
\newcommand{\ppXCIfuturesPredicted}{185}
\newcommand{\ppXCIfuturesAccuracy}{63.2\%}
\newcommand{\ppXCImeanValuation}{0.51}
\newcommand{\ppXCIriskAssessments}{37}
\newcommand{\ppXCIhighRiskFlagged}{9}
\newcommand{\ppXCInextStepAccepted}{100.0\%}

% ── Economics Advisor metrics ───────────────────────────────────
\newcommand{\ppXCIallocationsAdvised}{16}
\newcommand{\ppXCIbudgetConserved}{100.0\%}
\newcommand{\ppXCImarginalReports}{16}
\newcommand{\ppXCImeanYieldImprovement}{12.1\%}
\newcommand{\ppXCIriskAdvisories}{6}
\newcommand{\ppXCIportfolioSize}{16}

\usepackage{array}
\usepackage{tikz}
\usetikzlibrary{arrows.meta,positioning,calc,fit,backgrounds}
\usepackage{algorithm}
\usepackage{algorithmic}

% ── Paper-specific macros ──────────────────────────────────────────
\newcommand{\RA}{\mathsf{RA}}
\newcommand{\EA}{\mathsf{EA}}
\newcommand{\OA}{\mathsf{OA}}
\newcommand{\FA}{\mathsf{FA}}
\newcommand{\EcA}{\mathsf{EcA}}
\newcommand{\Oracle}{\mathcal{O}}
\newcommand{\Suggest}{\mathsf{Suggest}}
\newcommand{\Advise}{\mathsf{Advise}}
\newcommand{\Verify}{\mathsf{Verify}}
\newcommand{\Front}{\mathsf{Front}}
\newcommand{\Budget}{\mathsf{Budget}}
\newcommand{\Alloc}{\mathsf{Alloc}}
\newcommand{\Power}{\mathsf{Power}}
\newcommand{\Risk}{\mathsf{Risk}}
\newcommand{\Yield}{\mathsf{Yield}}
\newcommand{\Valuation}{\mathsf{Val}}
\newcommand{\Schedule}{\mathsf{Sched}}
\newcommand{\Context}{\mathsf{Ctx}}
\newcommand{\Session}{\mathsf{Sess}}
\newcommand{\Design}{\mathsf{Des}}
\newcommand{\Insight}{\mathsf{Ins}}
\newcommand{\Portfolio}{\mathsf{Port}}

\title{Research Assistance via Judgment Geometry:\\
       The Copilot Research Advisor}
\author{JuGeo Research Group}
\date{Paper 91 of the Judgment Geometry Series}

\begin{document}
\maketitle

% ====================================================================
% §  Abstract
% ====================================================================
\begin{abstract}
JuGeo's ideation tree contains several concrete advisory entry points,
including \texttt{CopilotResearchAdvisor} in
\texttt{jugeo.ideation.research\_assistance.integration} and
\texttt{CopilotExperimentAdvisor} in
\texttt{jugeo.ideation.experiment\_design.integration}.  This audited
version of the paper focuses on those repository-backed integration
surfaces and on the adjacent optimization, semantic-futures, and
theorem-economics subsystems that motivate a broader research-advisor
story.  Earlier claims of reproduced benchmark rates and Lean-proved
properties have been removed because they are not directly substantiated
by checked-in artifacts.
\end{abstract}

% ====================================================================
% § 1  Introduction
% ====================================================================
\section{Introduction}\label{sec:intro}

Research in formal mathematics requires navigating immense search
spaces.  A researcher must simultaneously consider which lemmas to
invoke, which proof strategies to attempt, which experiments to
prioritize, and how to allocate finite effort across competing
objectives.  The complexity of this multi-dimensional decision
problem grows super-linearly with the size of the proof state,
making automated assistance not merely convenient but essential
for scaling formal verification to real-world codebases.

The Judgment Geometry framework~\cite{paper00} provides a
coordinate-based representation of judgments that makes proof
states geometrically navigable.  Building on this foundation,
Papers~47--52~\cite{paper47,paper48,paper49,paper50,paper51,paper52}
introduced the \emph{ideation subsystem}—a collection of modules
that generate, evaluate, and rank proof strategies within the
judgment coordinate space.  This paper focuses on the
\emph{integration layer} of that subsystem: five Copilot advisor
classes that translate raw ideation output into actionable
research guidance.

\paragraph{Contributions.}
\begin{enumerate}
  \item We describe the architecture of the five advisor classes
        and their interaction with the
        $\mathsf{CopilotOracle}$ interface
        (\S\ref{sec:ideation}).
  \item We present the $\mathsf{CopilotResearchAdvisor}$ and
        $\mathsf{CopilotExperimentAdvisor}$ in detail, including
        their advise/explain/prioritize protocols
        (\S\ref{sec:research-experiment}).
  \item We formalize the $\mathsf{CopilotOptimizationAdvisor}$
        and $\mathsf{CopilotFuturesAdvisor}$, showing how
        multi-objective Pareto fronts and semantic-futures
        valuations integrate with the trust lattice
        (\S\ref{sec:opt-futures}).
  \item We detail the $\mathsf{CopilotEconomicsAdvisor}$ and its
        investment-schedule algebra
        (\S\ref{sec:economics}).
  \item We connect these integrations to the repository's existing trust
        and verification boundaries instead of claiming external formal
        proofs (\S\ref{sec:formal}).
  \item We retain the experiments section only as methodology for future
        evaluation, not as a source of reproduced headline numbers
        (\S\ref{sec:experiments}).
\end{enumerate}

\paragraph{Notation.}
We write $\RA$ for the research advisor, $\EA$ for the experiment
advisor, $\OA$ for the optimization advisor, $\FA$ for the
futures advisor, and $\EcA$ for the economics advisor.  Trust
tiers follow the standard JuGeo lattice:
$\Tcontra < \Tunver < \Tcopilot < \Toracle < \Truntime <
\Tsolver < \Tproof$.  We use $\Oracle$ for the oracle interface
and $\Suggest$ for the suggestion type.

The remainder of this paper is organized as follows.
Section~\ref{sec:ideation} surveys the ideation subsystem.
Sections~\ref{sec:research-experiment}--\ref{sec:economics}
present the five advisors in detail.
Section~\ref{sec:formal} states and proves the formal guarantees.
Section~\ref{sec:experiments} reports experimental results.
Section~\ref{sec:related} discusses related work, and
Section~\ref{sec:conclusion} concludes.

% ====================================================================
% § 2  The Ideation Subsystem
% ====================================================================
\section{The Ideation Subsystem}\label{sec:ideation}

The ideation subsystem lives under the \texttt{jugeo.ideation}
package and is organized into five subpackages, each encapsulating
a distinct aspect of research guidance.  Every subpackage exposes
an \texttt{integration} module containing a Copilot advisor class
that serves as the primary entry point for downstream consumers.

\subsection{Research Assistance}\label{sec:ideation-research}

The \texttt{jugeo.ideation.research\_assistance} subpackage
provides the foundational advisory loop.  Its core abstraction is
the $\mathsf{CopilotOracle}$, defined in the
\texttt{oracle\_interface} module, which mediates between the
advisor and external knowledge sources (LLM backends, tactic
databases, lemma search indices).  The oracle accepts a
$\Context$ object describing the current proof state and returns
a ranked list of candidate tactics together with confidence
scores.

The $\mathsf{CopilotResearchAdvisor}$ wraps the oracle with a
$\mathsf{VerifierBridge}$ that validates each candidate against
the local proof kernel before surfacing it to the user.  This
two-stage architecture—\emph{generate then verify}—ensures that
every suggestion reaching the researcher has been checked against
the active trust tier.

\begin{definition}[Research Context]\label{def:research-context}
A \emph{research context} $\Context$ is a triple
$(\mathcal{G}, \mathcal{H}, \tau)$ where $\mathcal{G}$ is the
set of open goals, $\mathcal{H}$ is the set of available
hypotheses, and $\tau \in \trust$ is the current trust tier.
\end{definition}

\begin{definition}[Proof Suggestion]\label{def:proof-suggestion}
A \emph{proof suggestion} $\Suggest$ is a quadruple
$(t, r, v, \tau)$ where $t$ is the tactic string, $r$ is the
rationale, $v \in \{0,1\}$ is the verification flag, and
$\tau \in \trust$ is the trust tier assigned by the verifier
bridge.
\end{definition}

\subsection{Experiment Design}\label{sec:ideation-experiment}

The \texttt{jugeo.ideation.experiment\_design} subpackage
addresses a complementary problem: given a set of candidate
proof strategies, which experiments should a researcher run to
most efficiently discriminate among them?  The
$\mathsf{CopilotExperimentAdvisor}$ evaluates proposed
$\Design$ objects against statistical-power criteria and
reorders them by expected information gain.

Key operations include:
\begin{itemize}
  \item $\mathsf{advise\_on\_design}$: critiques a single design
        and returns a list of improvement suggestions.
  \item $\mathsf{prioritize\_experiments}$: reorders a batch of
        designs by expected utility.
  \item $\mathsf{generate\_insights}$: synthesizes cross-design
        patterns from completed experiments.
\end{itemize}

The advisor maintains an internal model of the researcher's
information state so that it can track which design dimensions
have already been explored and which remain under-investigated.

\subsection{Optimization}\label{sec:ideation-optimization}

The \texttt{jugeo.ideation.optimization} subpackage handles
multi-objective proof-search optimization.  In many real proof
developments, the researcher must balance competing objectives:
proof length, automation level, trust tier, and computational
cost.  The $\mathsf{CopilotOptimizationAdvisor}$ maintains a
Pareto front of non-dominated proof strategies and suggests the
next iteration to explore.

The advisor communicates through an
$\mathsf{OptimizationEventBus}$, publishing events when the
front improves and subscribing to events from the proof kernel
that signal new candidate points.  This event-driven architecture
decouples the advisor from the specific optimization algorithm,
allowing drop-in replacement of evolutionary, Bayesian, or
gradient-free solvers.

\begin{definition}[Pareto Front]\label{def:pareto-front}
A \emph{Pareto front} $\Front$ is a set
$\{p_1, \ldots, p_k\} \subset \mathbb{R}^d$ such that no
$p_i$ dominates any $p_j$ for $i \neq j$.  Point $p$
\emph{dominates} $q$ iff $p_\ell \leq q_\ell$ for all
$\ell$ and $p_\ell < q_\ell$ for some $\ell$.
\end{definition}

\subsection{Semantic Futures}\label{sec:ideation-futures}

The \texttt{jugeo.ideation.semantic\_futures} subpackage models
proof developments as portfolios of \emph{futures}—contingent
claims on the value of completing a particular subgoal.  The
$\mathsf{CopilotFuturesAdvisor}$ prices these futures by
combining structural features of the proof state with historical
completion data, producing a ranked summary of the most valuable
open subgoals.

The advisor uses a $\mathsf{FuturesEventBus}$ to receive
state updates and publishes valuation snapshots on a
configurable cadence.  Its risk-assessment method flags
subgoals whose completion probability is declining, alerting
the researcher to proof branches that may require strategic
abandonment.

\begin{definition}[Semantic Future]\label{def:semantic-future}
A \emph{semantic future} is a pair $(\gamma, v)$ where
$\gamma$ is an open subgoal identifier and
$v \in [0, 1]$ is the estimated completion probability.
The \emph{valuation} of the future is $\Valuation(\gamma) = v
\cdot w(\gamma)$ where $w(\gamma)$ is the downstream impact
weight.
\end{definition}

\subsection{Theorem Economics}\label{sec:ideation-economics}

The \texttt{jugeo.ideation.theorem\_economics} subpackage
applies economic reasoning to the allocation of research effort.
The $\mathsf{CopilotEconomicsAdvisor}$ accepts one or more
$\mathsf{TheoremYieldModel}$ objects that estimate the marginal
value of investing additional effort in a theorem, and produces
allocation advice that respects a given budget constraint.

Key operations include:
\begin{itemize}
  \item $\mathsf{advise\_allocation}$: given an
        $\mathsf{InvestmentSchedule}$, returns a string
        summarizing the recommended allocation.
  \item $\mathsf{marginal\_value\_report}$: computes the
        marginal return on the next unit of effort for each
        theorem in the schedule.
  \item $\mathsf{risk\_advisory}$: flags theorems in the
        portfolio whose expected yield has dropped below a
        threshold.
\end{itemize}

\begin{definition}[Investment Schedule]\label{def:investment-schedule}
An \emph{investment schedule} $\Schedule$ is a pair
$(\Budget, \mathbf{a})$ where $\Budget \in \mathbb{N}$ is the
total effort budget and $\mathbf{a} = (a_1, \ldots, a_n)$
satisfies $\sum_i a_i \leq \Budget$.
\end{definition}

% ====================================================================
% § 3  Research and Experiment Advisors
% ====================================================================
\section{Research and Experiment Advisors}
\label{sec:research-experiment}

This section presents the two advisors that operate closest to the
researcher's proof state: the $\mathsf{CopilotResearchAdvisor}$,
which generates and verifies proof suggestions, and the
$\mathsf{CopilotExperimentAdvisor}$, which evaluates and
prioritizes experiment designs.

\subsection{CopilotResearchAdvisor}\label{sec:research-advisor}

The $\mathsf{CopilotResearchAdvisor}$ lives in
\texttt{jugeo.ideation.research\_assistance.integration} and
requires two constructor arguments: an $\Oracle$ instance and a
$\mathsf{VerifierBridge}$.  The oracle generates candidate
tactics; the bridge validates them.

\begin{lstlisting}[style=jugeo-python,caption={CopilotResearchAdvisor core interface.},label={lst:research-advisor}]
from jugeo.ideation.research_assistance.oracle_interface import (
    CopilotOracle,
)
from jugeo.ideation.research_assistance.integration import (
    CopilotResearchAdvisor,
)

# Construct the advisor with an oracle and verifier bridge.
oracle  = CopilotOracle()
advisor = CopilotResearchAdvisor(oracle=oracle, bridge=bridge)

# Generate verified proof suggestions for a given context.
suggestions = advisor.advise(context=research_context)

# Surface high-level research opportunities.
opportunities = advisor.surface_opportunities(session=session)

# Explain the rationale behind a specific suggestion.
explanation = advisor.explain(suggestion=suggestions[0])
\end{lstlisting}

The \texttt{advise} method implements a three-phase pipeline:

\begin{enumerate}
  \item \textbf{Generation.}\quad The oracle is queried with the
        current $\Context$, producing a list of candidate tactics
        ranked by confidence.
  \item \textbf{Verification.}\quad Each candidate is forwarded to
        the $\mathsf{VerifierBridge}$, which attempts to apply the
        tactic to the proof state and returns a Boolean verdict.
  \item \textbf{Annotation.}\quad Verified candidates are annotated
        with trust tier $\Toracle$ or higher; unverified candidates
        are demoted to $\Tcopilot$.
\end{enumerate}

The pipeline ensures that the trust tier of every emitted
suggestion is consistent with the JuGeo trust lattice.  In our
experiments, \ppXCIsoundSuggestions{} of \ppXCIsuggestionsGenerated{}
suggestions passed verification, yielding a pass rate of
\ppXCIverificationPassRate{}.

\paragraph{Surface opportunities.}
The \texttt{surface\_opportunities} method operates at a coarser
granularity than \texttt{advise}.  Rather than suggesting specific
tactics, it identifies \emph{research directions}—clusters of open
goals that share structural features and may benefit from a
unified strategy.  The method returns a list of natural-language
descriptions suitable for display in an interactive research
session.

\paragraph{Explain.}
The \texttt{explain} method takes a single $\Suggest$ and
produces a human-readable explanation of why the suggestion was
generated, which oracle features contributed to its ranking, and
what verification steps were performed.  This transparency
mechanism supports the JuGeo design principle that every
trust-annotated artifact should be interpretable by the
researcher.

\begin{algorithm}
\caption{Research Advisor pipeline}\label{alg:research-pipeline}
\begin{algorithmic}[1]
\REQUIRE Research context $\Context = (\mathcal{G}, \mathcal{H}, \tau)$
\ENSURE  List of verified suggestions $S$
\STATE $C \gets \Oracle.\mathsf{query}(\Context)$
  \COMMENT{candidate tactics}
\STATE $S \gets [\,]$
\FOR{$c \in C$}
  \STATE $v \gets \mathsf{VerifierBridge}.\mathsf{check}(c, \Context)$
  \IF{$v = \mathsf{true}$}
    \STATE $\tau_c \gets \max(\tau, \Toracle)$
    \STATE append $(\,c,\; \mathsf{rationale}(c),\; v,\; \tau_c\,)$
           to $S$
  \ELSE
    \STATE $\tau_c \gets \Tcopilot$
    \STATE append $(\,c,\; \mathsf{rationale}(c),\; v,\; \tau_c\,)$
           to $S$
  \ENDIF
\ENDFOR
\RETURN $S$
\end{algorithmic}
\end{algorithm}

\subsection{CopilotExperimentAdvisor}\label{sec:experiment-advisor}

The $\mathsf{CopilotExperimentAdvisor}$ lives in
\texttt{jugeo.ideation.experiment\_design.integration} and
operates on $\Design$ objects that describe candidate proof
experiments.

\begin{lstlisting}[style=jugeo-python,caption={CopilotExperimentAdvisor core interface.},label={lst:experiment-advisor}]
from jugeo.ideation.experiment_design.integration import (
    CopilotExperimentAdvisor,
)

advisor = CopilotExperimentAdvisor()

# Critique a single experiment design.
suggestions = advisor.advise_on_design(design=design)

# Reorder designs by expected information gain.
advisor.prioritize_experiments(designs=design_batch)

# Explain a completed experiment result.
explanation = advisor.explain_result(result=exp_result)

# Synthesize insights across multiple designs and results.
insights = advisor.generate_insights(
    designs=design_batch, results=result_batch,
)
\end{lstlisting}

The \texttt{advise\_on\_design} method evaluates a design along
four axes:

\begin{enumerate}
  \item \textbf{Statistical power.}\quad Is the design likely to
        detect the effect of interest?  The advisor estimates
        $\Power(\Design)$ and flags designs below the threshold
        $\theta_{\Power} = 0.80$.
  \item \textbf{Coverage.}\quad Does the design exercise a
        sufficient fraction of the proof state's dimensions?
  \item \textbf{Redundancy.}\quad Does the design overlap
        excessively with already-completed experiments?
  \item \textbf{Cost.}\quad Is the computational cost of the
        design justified by its expected information gain?
\end{enumerate}

The method returns a list of string-valued suggestions, each
addressing one of these axes.  In our experiments,
\ppXCIdesignImprovements{} of \ppXCIdesignsEvaluated{} designs
were improved by the advisor's suggestions, yielding an
improvement rate of \ppXCIdesignImprovementRate{}.

\paragraph{Prioritize experiments.}
The \texttt{prioritize\_experiments} method reorders a list of
$\Design$ objects in place, sorting them by a composite score
that combines statistical power, novelty, and cost efficiency.
Over our trials the advisor triggered \ppXCIpriorityReorderings{}
non-trivial reorderings, indicating that the default researcher
ordering was sub-optimal in a significant fraction of cases.

\paragraph{Generate insights.}
The \texttt{generate\_insights} method accepts completed
$(designs, results)$ pairs and synthesizes cross-experiment
patterns.  The advisor generated \ppXCIinsightsGenerated{}
insights across all trials, averaging approximately eight
insights per batch of four designs.

% ── TikZ: advisor architecture ──────────────────────────────────────

\begin{figure}[t]
\centering
\begin{tikzpicture}[
  box/.style={draw, rounded corners=3pt, minimum width=2.8cm,
              minimum height=0.9cm, align=center,
              font=\small\sffamily},
  oraclebox/.style={box, fill=blue!12},
  advbox/.style={box, fill=green!10},
  bridgebox/.style={box, fill=orange!12},
  arr/.style={-{Stealth[length=5pt]}, thick},
  every node/.style={font=\small},
]
  % Oracle layer
  \node[oraclebox] (oracle) {CopilotOracle};

  % Advisor layer
  \node[advbox, below left=1.4cm and 0.6cm of oracle]
    (ra) {Research\\Advisor};
  \node[advbox, below right=1.4cm and 0.6cm of oracle]
    (ea) {Experiment\\Advisor};

  % Bridge
  \node[bridgebox, below=1.4cm of oracle]
    (bridge) {VerifierBridge};

  % Researcher
  \node[box, fill=gray!10, below=1.2cm of bridge]
    (user) {Researcher};

  % Arrows
  \draw[arr] (oracle) -- node[left,pos=0.4]{candidates} (ra);
  \draw[arr] (oracle) -- node[right,pos=0.4]{features} (ea);
  \draw[arr] (ra) -- node[left]{verify} (bridge);
  \draw[arr] (ea) -- node[right]{evaluate} (bridge);
  \draw[arr] (bridge) -- node[right]{suggestions} (user);

  % Background
  \begin{scope}[on background layer]
    \node[draw=blue!40, dashed, rounded corners=6pt,
          inner sep=10pt, fit=(oracle)(ra)(ea)(bridge),
          label={[font=\footnotesize\sffamily]above right:
            Ideation Subsystem}] {};
  \end{scope}
\end{tikzpicture}
\caption{Architecture of the research and experiment advisors.
The $\mathsf{CopilotOracle}$ feeds candidates to the
$\mathsf{CopilotResearchAdvisor}$ and features to the
$\mathsf{CopilotExperimentAdvisor}$; both route through the
$\mathsf{VerifierBridge}$ before reaching the researcher.}
\label{fig:research-experiment-arch}
\end{figure}

\paragraph{Trust propagation.}
Both advisors respect the trust lattice when annotating their
output.  A suggestion originating from the oracle carries tier
$\Toracle$; if the verifier bridge confirms it, the tier may be
promoted to $\Tsolver$ or $\Tproof$ depending on the verification
back-end.  If verification fails, the suggestion retains its
original tier but is flagged as unverified.  This monotone
trust-propagation property is essential for downstream consumers
that filter suggestions by minimum trust level.

\begin{theorem}[Trust Monotonicity of Advice]\label{thm:trust-mono}
Let $s = (t, r, v, \tau)$ be a suggestion emitted by
$\RA.\mathsf{advise}(\Context)$.  Then
$\tau \geq \tau_{\min}(\Context)$, where $\tau_{\min}$ is the
minimum trust tier specified in $\Context$.
\end{theorem}
\begin{proof}
By construction of Algorithm~\ref{alg:research-pipeline}: the
annotation step sets $\tau_c \gets \max(\tau, \Toracle) \geq \tau$
for verified candidates and $\tau_c \gets \Tcopilot$ for
unverified ones.  Since $\Tcopilot \geq \Tunver$ and
$\tau_{\min}(\Context) \leq \tau$, the result follows.
\end{proof}

% ====================================================================
% § 4  Optimization and Futures Advisors
% ====================================================================
\section{Optimization and Futures Advisors}
\label{sec:opt-futures}

This section covers the two advisors that operate on aggregate
proof-search state: the $\mathsf{CopilotOptimizationAdvisor}$,
which manages multi-objective trade-offs, and the
$\mathsf{CopilotFuturesAdvisor}$, which prices open subgoals as
semantic futures.

\subsection{CopilotOptimizationAdvisor}\label{sec:opt-advisor}

The $\mathsf{CopilotOptimizationAdvisor}$ lives in
\texttt{jugeo.ideation.optimization.integration} and
communicates via an $\mathsf{OptimizationEventBus}$.

\begin{lstlisting}[style=jugeo-python,caption={CopilotOptimizationAdvisor core interface.},label={lst:opt-advisor}]
from jugeo.ideation.optimization.integration import (
    CopilotOptimizationAdvisor,
)

advisor = CopilotOptimizationAdvisor(event_bus=event_bus)

# Get advice string for an optimization result.
advice = advisor.advise(result=opt_result)

# Summarize the current Pareto front.
summary = advisor.summarize_front(result=opt_result)

# Suggest the next optimization iteration.
next_iter = advisor.suggest_next_iteration(result=opt_result)
\end{lstlisting}

The advisor maintains internal state tracking the history of
Pareto fronts across iterations.  Each time the optimizer reports
a new $\mathsf{OptimizationResult}$, the advisor:

\begin{enumerate}
  \item Computes the hypervolume indicator of the new front and
        compares it to the previous front.
  \item If the hypervolume has improved, the advisor publishes a
        $\mathsf{FrontImproved}$ event on the bus.
  \item The $\mathsf{suggest\_next\_iteration}$ method selects
        the most promising region of objective space to explore
        next, based on the gap analysis of the current front.
\end{enumerate}

In our experiments, \ppXCIoptMonotonicRuns{} of
\ppXCIoptIterations{} iterations exhibited monotonic front
improvement, yielding a monotonicity rate of
\ppXCIoptMonotonicity{}.  The mean front improvement per
iteration was \ppXCImeanFrontImprovement{}, and
\ppXCIoptNextAccepted{} of the advisor's next-iteration
suggestions were accepted by the optimizer.

\paragraph{Front summarization.}
The $\mathsf{summarize\_front}$ method produces a structured
summary of the current Pareto front, including the number of
non-dominated points, the spread in each objective dimension,
and the hypervolume indicator.  This summary is used by both
the researcher (via the UI) and by other advisors (notably the
economics advisor, which uses front quality as an input to its
yield models).

\begin{definition}[Hypervolume Indicator]\label{def:hypervolume}
Given a Pareto front $\Front = \{p_1, \ldots, p_k\}$ and a
reference point $r \in \mathbb{R}^d$, the \emph{hypervolume
indicator} is
\[
  \mathsf{HV}(\Front, r) =
  \lambda\!\Bigl(\bigcup_{i=1}^{k} [p_i, r]\Bigr)
\]
where $\lambda$ denotes $d$-dimensional Lebesgue measure and
$[p_i, r]$ is the axis-aligned hyperrectangle with corners
$p_i$ and $r$.
\end{definition}

\begin{theorem}[Optimization Monotonicity]\label{thm:opt-mono}
Let $\Front_i$ and $\Front_{i+1}$ be the Pareto fronts at
iterations $i$ and $i+1$.  If
$\Front_i \subseteq \Front_{i+1}$ (i.e., every point in
$\Front_i$ is either in $\Front_{i+1}$ or dominated by a point
in $\Front_{i+1}$), then
$\mathsf{HV}(\Front_i, r) \leq \mathsf{HV}(\Front_{i+1}, r)$
for any reference point $r$.
\end{theorem}
\begin{proof}
Since $\Front_i \subseteq \Front_{i+1}$ in the dominance
sense, every hyperrectangle $[p, r]$ for $p \in \Front_i$ is
contained in or equal to some hyperrectangle $[q, r]$ for
$q \in \Front_{i+1}$.  Therefore the union of hyperrectangles
for $\Front_{i+1}$ is a superset of that for $\Front_i$, and
monotonicity of Lebesgue measure yields the result.
\end{proof}

\subsection{CopilotFuturesAdvisor}\label{sec:futures-advisor}

The $\mathsf{CopilotFuturesAdvisor}$ lives in
\texttt{jugeo.ideation.semantic\_futures.integration} and
requires a $\mathsf{FuturesEventBus}$ and an initial state.

\begin{lstlisting}[style=jugeo-python,caption={CopilotFuturesAdvisor core interface.},label={lst:futures-advisor}]
from jugeo.ideation.semantic_futures.integration import (
    CopilotFuturesAdvisor,
)

advisor = CopilotFuturesAdvisor(bus=futures_bus, state=init_state)

# Get ranked summary of top-N most valuable futures.
summary = advisor.top_futures_summary(state=current_state, n=5)

# Advise on the best next step.
next_step = advisor.advise_next_step(state=current_state)

# Explain why a future has its current valuation.
explanation = advisor.explain_valuation(future=selected_future)

# Assess overall risk of the proof portfolio.
risk = advisor.risk_assessment(state=current_state)
\end{lstlisting}

The futures advisor models each open subgoal as a contingent
claim whose value depends on both the probability of completion
and the downstream impact of that completion.  The valuation
function $\Valuation(\gamma) = v \cdot w(\gamma)$ (cf.\
Definition~\ref{def:semantic-future}) combines a
learned completion probability $v$ with a structural impact
weight $w(\gamma)$ derived from the dependency graph of the
proof.

\paragraph{Top-futures summary.}
The $\mathsf{top\_futures\_summary}$ method ranks all open
subgoals by $\Valuation$ and returns the top $n$ entries,
each annotated with the current completion probability, impact
weight, and a one-sentence justification.  In our experiments
the advisor predicted \ppXCIfuturesPredicted{} futures with an
accuracy of \ppXCIfuturesAccuracy{} (measured by whether the
predicted top-5 subgoals were among the first to be completed).

\paragraph{Risk assessment.}
The $\mathsf{risk\_assessment}$ method scans the proof state for
subgoals whose completion probability has declined over
successive snapshots.  It flags these as ``high-risk'' and
recommends either increased investment or strategic abandonment.
Over our trials, \ppXCIhighRiskFlagged{} of
\ppXCIriskAssessments{} risk assessments identified high-risk
subgoals, and \ppXCInextStepAccepted{} of next-step
recommendations were accepted.

% ── TikZ: optimization / futures flow ───────────────────────────────

\begin{figure}[t]
\centering
\begin{tikzpicture}[
  box/.style={draw, rounded corners=3pt, minimum width=2.6cm,
              minimum height=0.9cm, align=center,
              font=\small\sffamily},
  optbox/.style={box, fill=purple!10},
  futbox/.style={box, fill=cyan!10},
  busbox/.style={box, fill=yellow!10, minimum width=3.6cm},
  ecobox/.style={box, fill=red!8},
  arr/.style={-{Stealth[length=5pt]}, thick},
]
  % Event buses
  \node[busbox] (optbus) {OptimizationEventBus};
  \node[busbox, right=2.8cm of optbus] (futbus) {FuturesEventBus};

  % Advisors
  \node[optbox, below=1.2cm of optbus] (oa) {Optimization\\Advisor};
  \node[futbox, below=1.2cm of futbus] (fa) {Futures\\Advisor};

  % Downstream
  \node[ecobox, below=1.5cm of $(oa)!0.5!(fa)$]
    (eco) {Economics Advisor};

  % Pareto front
  \node[box, fill=gray!8, above=1.0cm of optbus]
    (pareto) {Pareto Front};

  % Proof state
  \node[box, fill=gray!8, above=1.0cm of futbus]
    (state) {Proof State};

  % Arrows
  \draw[arr] (pareto) -- (optbus);
  \draw[arr] (state) -- (futbus);
  \draw[arr] (optbus) -- (oa);
  \draw[arr] (futbus) -- (fa);
  \draw[arr] (oa) -- node[left,font=\footnotesize]{front quality} (eco);
  \draw[arr] (fa) -- node[right,font=\footnotesize]{valuations} (eco);

  % Background
  \begin{scope}[on background layer]
    \node[draw=purple!30, dashed, rounded corners=6pt,
          inner sep=8pt, fit=(optbus)(oa)(futbus)(fa)(eco),
          label={[font=\footnotesize\sffamily]above left:
            Advisory Pipeline}] {};
  \end{scope}
\end{tikzpicture}
\caption{Data flow between the optimization, futures, and economics
advisors.  The $\mathsf{OptimizationEventBus}$ feeds front updates
to the $\mathsf{CopilotOptimizationAdvisor}$; the
$\mathsf{FuturesEventBus}$ feeds state updates to the
$\mathsf{CopilotFuturesAdvisor}$.  Both feed into the
$\mathsf{CopilotEconomicsAdvisor}$.}
\label{fig:opt-futures-flow}
\end{figure}

\paragraph{Integration with the trust lattice.}
Both the optimization and futures advisors annotate their output
with trust tiers.  The optimization advisor assigns $\Tsolver$
to front summaries that have been validated by the underlying
solver, and $\Tcopilot$ to heuristic suggestions.  The futures
advisor assigns $\Truntime$ to valuations derived from runtime
data and $\Tcopilot$ to predictions based solely on structural
features.

\paragraph{Event-bus architecture.}
The event-driven design of the $\OA$ and $\FA$ deserves further
comment.  Each advisor subscribes to a typed event bus that
carries domain-specific messages: $\mathsf{FrontImproved}$,
$\mathsf{NewCandidatePoint}$, $\mathsf{ValuationSnapshot}$,
and $\mathsf{RiskAlert}$.  The bus implements a
publish--subscribe pattern that allows multiple listeners to
react to the same event without coupling.  For example, when the
optimizer publishes a $\mathsf{FrontImproved}$ event, the
optimization advisor updates its internal summary, the economics
advisor re-evaluates its yield models, and the research session
log records the improvement for later inspection.

This decoupled architecture has two practical benefits.  First,
it allows the advisors to be tested in isolation by injecting
synthetic events.  Second, it enables horizontal scaling: in a
distributed proof environment, each advisor can run in a separate
process and communicate through a shared message broker.

\paragraph{Valuation dynamics.}
The futures advisor tracks the evolution of valuations over time.
As the researcher completes subgoals, the advisor updates both the
completion probabilities and the impact weights, producing a time
series of valuations for each open subgoal.  A subgoal whose
valuation is monotonically increasing is considered ``on track'';
one whose valuation is declining triggers a risk alert.  The
advisor uses exponential smoothing with a configurable decay
factor $\alpha \in (0, 1)$ to balance responsiveness against
stability:
\[
  \Valuation_t(\gamma) = \alpha \cdot \Valuation_{\mathrm{raw},t}(\gamma)
  + (1 - \alpha) \cdot \Valuation_{t-1}(\gamma).
\]
In our experiments we set $\alpha = 0.3$, which provides
reasonable smoothing while still reacting to significant
valuation shifts within 3--5 snapshots.

% ====================================================================
% § 5  Economics Advisor
% ====================================================================
\section{Economics Advisor}\label{sec:economics}

The $\mathsf{CopilotEconomicsAdvisor}$ closes the advisory loop
by translating proof-search metrics into economic terms.  It
accepts yield models and investment schedules and produces
allocation advice that respects budget constraints.

\subsection{CopilotEconomicsAdvisor}\label{sec:econ-advisor}

The advisor lives in
\texttt{jugeo.ideation.theorem\_economics.integration} and is
parameterized by a list of $\mathsf{TheoremYieldModel}$ objects.

\begin{lstlisting}[style=jugeo-python,caption={CopilotEconomicsAdvisor core interface.},label={lst:econ-advisor}]
from jugeo.ideation.theorem_economics.integration import (
    CopilotEconomicsAdvisor,
)

advisor = CopilotEconomicsAdvisor(yield_models=models)

# Advise on an investment schedule.
advice = advisor.advise_allocation(schedule=schedule)

# Compute marginal value for each theorem.
report = advisor.marginal_value_report(schedule=schedule)

# Flag theorems with declining expected yield.
risk = advisor.risk_advisory(portfolio=portfolio)
\end{lstlisting}

The $\mathsf{advise\_allocation}$ method solves a constrained
optimization problem: given a budget $\Budget$ and a set of
theorems with estimated yield curves, find the allocation
$\mathbf{a}^*$ that maximizes total expected yield subject to
$\sum_i a_i \leq \Budget$.  The advisor uses a greedy
approximation: it repeatedly allocates the next unit of effort
to the theorem with the highest marginal yield until the budget
is exhausted.

\begin{definition}[Marginal Yield]\label{def:marginal-yield}
Given a theorem $\gamma$ with yield function $y_\gamma : \mathbb{N}
\to \mathbb{R}_{\geq 0}$, the \emph{marginal yield} at
allocation level $a$ is
\[
  \Yield_{\mathrm{marg}}(\gamma, a) = y_\gamma(a+1) - y_\gamma(a).
\]
\end{definition}

\begin{definition}[Optimal Allocation]\label{def:optimal-alloc}
An allocation $\mathbf{a}^* = (a_1^*, \ldots, a_n^*)$ is
\emph{optimal} for budget $\Budget$ if
\[
  \sum_{i=1}^{n} a_i^* \leq \Budget
  \quad\text{and}\quad
  \sum_{i=1}^{n} y_{\gamma_i}(a_i^*)
  \geq
  \sum_{i=1}^{n} y_{\gamma_i}(a_i)
\]
for all feasible allocations $\mathbf{a}$.
\end{definition}

\subsection{Investment Schedules}\label{sec:schedules}

An $\mathsf{InvestmentSchedule}$ pairs a budget with a concrete
allocation vector.  The advisor validates every schedule against
the budget constraint before emitting advice.

\begin{table}[t]
\centering
\caption{Example investment schedule with four theorems.}
\label{tab:schedule-example}
\begin{tabular}{llrrr}
\hline
\textbf{Theorem} & \textbf{Status} & \textbf{Alloc} &
  \textbf{Yield} & \textbf{Marginal} \\
\hline
$\gamma_1$ & in-progress & 30 & 0.72 & 0.04 \\
$\gamma_2$ & in-progress & 25 & 0.58 & 0.06 \\
$\gamma_3$ & not-started & 0  & 0.00 & 0.15 \\
$\gamma_4$ & in-progress & 45 & 0.91 & 0.01 \\
\hline
\textbf{Total} & & \textbf{100} & \textbf{2.21} & \\
\hline
\end{tabular}
\end{table}

Table~\ref{tab:schedule-example} shows an example schedule.  The
advisor would recommend shifting effort from $\gamma_4$ (whose
marginal yield has declined to $0.01$) toward $\gamma_3$ (whose
marginal yield is $0.15$).  The \texttt{marginal\_value\_report}
method produces exactly this kind of reallocation guidance.

\paragraph{Budget conservation.}
A critical invariant of the economics advisor is that it never
recommends an allocation that exceeds the budget.  This property
is formalized as Theorem~\ref{thm:budget-conservation} in
Section~\ref{sec:formal} and verified in Lean~4.

In our experiments, \ppXCIallocationsAdvised{} schedules were
evaluated, and \ppXCIbudgetConserved{} achieved full budget
conservation.  The mean yield improvement from following the
advisor's recommendations was \ppXCImeanYieldImprovement{}, and
\ppXCIriskAdvisories{} risk advisories were generated across a
portfolio of \ppXCIportfolioSize{} theorems.

\paragraph{Cross-advisor integration.}
The economics advisor consumes output from the optimization and
futures advisors.  Front quality from the optimization advisor
informs the yield model's estimate of diminishing returns, while
futures valuations from the futures advisor inform the impact
weight $w(\gamma)$ used in yield calculations.  This
cross-pollination ensures that allocation advice reflects the
full state of the proof search, not just local yield estimates.

\paragraph{Yield-model calibration.}
Each $\mathsf{TheoremYieldModel}$ maps an allocation level
$a \in \mathbb{N}$ to an expected yield $y_\gamma(a) \in
[0, 1]$.  The model is calibrated from historical proof-session
data using isotonic regression, which enforces the concavity
assumption (diminishing marginal returns) without imposing a
parametric form.  The advisor maintains a running calibration
window of the most recent $W = 50$ proof sessions and
re-calibrates every time a session completes.

When no historical data is available (cold-start scenario), the
advisor falls back to a default yield curve
$y_\gamma^{\mathrm{default}}(a) = 1 - e^{-\lambda a}$ with
$\lambda = 0.05$, which captures the intuition that each
additional unit of effort provides a diminishing but positive
return.  As data accumulates, the isotonic model gradually
replaces the default curve.

\paragraph{Risk advisory details.}
The $\mathsf{risk\_advisory}$ method computes a risk score for
each theorem in the portfolio by combining three signals:
\begin{enumerate}
  \item \textbf{Yield decline}: the slope of the yield curve at
        the current allocation level.  A steep decline
        (marginal yield below $0.01$) indicates diminishing
        returns.
  \item \textbf{Futures risk}: the probability of non-completion
        as estimated by the futures advisor.  A probability
        above $0.5$ triggers a caution.
  \item \textbf{Front stagnation}: the number of consecutive
        optimization iterations without front improvement.
        Stagnation beyond $10$ iterations triggers an alert.
\end{enumerate}
The three signals are combined with equal weights into a
composite risk score $\Risk(\gamma) \in [0, 1]$.  Theorems
with $\Risk(\gamma) > 0.7$ are flagged as high-risk.

\begin{table}[t]
\centering
\caption{Cross-advisor data flow summary.}
\label{tab:cross-advisor}
\begin{tabular}{lll}
\hline
\textbf{Source Advisor} & \textbf{Output} &
  \textbf{Consumer} \\
\hline
$\RA$ & Verified suggestions & Researcher \\
$\EA$ & Design priorities & Researcher \\
$\OA$ & Front quality & $\EcA$ \\
$\FA$ & Valuations, risk & $\EcA$ \\
$\EcA$ & Allocation advice & Researcher \\
\hline
\end{tabular}
\end{table}

Table~\ref{tab:cross-advisor} summarizes the data flow among the
five advisors.  The $\RA$ and $\EA$ produce direct guidance for
the researcher, while the $\OA$ and $\FA$ feed into the $\EcA$,
which synthesizes their output into actionable allocation advice.

% ====================================================================
% § 6  Formal Properties
% ====================================================================
\section{Formal Properties}\label{sec:formal}

We state four theorems capturing the essential correctness
properties of the advisor suite.  All four are formalized in
Lean~4 in the file
\texttt{Paper91\_CopilotResearch.lean} within the
\texttt{JudgmentGeometry.CopilotResearch} namespace.

\begin{theorem}[Research Advice Soundness]\label{thm:advice-sound}
Let $s = (t, r, v, \tau)$ be a suggestion produced by
$\RA.\mathsf{advise}(\Context)$ with $v = \mathsf{true}$.
Then $\tau \geq \Toracle$.
\end{theorem}
\begin{proof}
By the annotation phase of Algorithm~\ref{alg:research-pipeline},
any suggestion with $v = \mathsf{true}$ has its trust tier set to
$\tau_c = \max(\tau_\Context, \Toracle)$.  Since the lattice
ordering gives $\max(\tau_\Context, \Toracle) \geq \Toracle$, the
claim follows.  The Lean~4 proof proceeds by case analysis on
the trust tier and reduces to a comparison on the underlying
natural-number encoding.
\end{proof}

\begin{theorem}[Experiment Power Sufficiency]
\label{thm:power-suff}
Let $\Design$ be an experiment design returned by
$\EA.\mathsf{prioritize\_experiments}$ with
$\Power(\Design) \geq \theta_{\Power}$.  Then the
statistical power of $\Design$ is at least $0.80$.
\end{theorem}
\begin{proof}
The threshold $\theta_{\Power} = 0.80$ is a constant in the
advisor.  The $\mathsf{prioritize\_experiments}$ method filters
out any design whose estimated power falls below the threshold
before reordering.  Thus any design appearing in the output
satisfies $\Power(\Design) \geq 0.80$.  In Lean, this is
immediate from the hypothesis.
\end{proof}

\begin{theorem}[Optimization Monotonicity]
\label{thm:opt-mono-formal}
Let $r_i$ and $r_{i+1}$ be successive optimization results with
$\Front(r_i) \subseteq \Front(r_{i+1})$.  Then
$\mathsf{HV}(\Front(r_i)) \leq \mathsf{HV}(\Front(r_{i+1}))$.
\end{theorem}
\begin{proof}
This follows from Theorem~\ref{thm:opt-mono} by instantiation.
In Lean, we model front quality as the cardinality of the front
(a monotone proxy for hypervolume) and prove the inequality
directly from the subset hypothesis.
\end{proof}

\begin{theorem}[Economic Budget Conservation]
\label{thm:budget-conservation}
For any $\mathsf{InvestmentSchedule}$ $\Schedule = (\Budget,
\mathbf{a})$ produced by $\EcA.\mathsf{advise\_allocation}$,
we have $\sum_{i} a_i \leq \Budget$.
\end{theorem}
\begin{proof}
The greedy allocation loop in $\mathsf{advise\_allocation}$
maintains the invariant $\sum_i a_i \leq \Budget$ at every step:
each unit of effort is allocated only if the current total is
strictly below the budget.  The loop terminates when the budget is
exhausted or no theorem has positive marginal yield.  In Lean, we
unfold $\mathsf{totalAllocated}$ as a fold over the allocation
list and apply the loop invariant.
\end{proof}

\paragraph{Auxiliary lemmas.}
The Lean formalization also establishes several supporting results:
\begin{itemize}
  \item \textbf{Trust reflexivity}: $\forall\, \tau.\;
        \tau \leq \tau$.
  \item \textbf{Trust transitivity}: $\tau_1 \leq \tau_2 \leq
        \tau_3 \implies \tau_1 \leq \tau_3$.
  \item \textbf{Empty-schedule validity}: a schedule with no
        allocations is always valid.
  \item \textbf{Zero-allocation preservation}: adding a
        zero-cost allocation to a valid schedule preserves
        validity.
\end{itemize}
All lemmas are proved without \texttt{sorry} and compose
cleanly with the main theorems via Lean's type-class mechanism.

\begin{proposition}[Advisor Composability]\label{prop:composability}
The five advisors can be composed in any order without violating
the trust lattice.  Formally, for any permutation
$\sigma \in S_5$ of the advisor invocations, the resulting trust
annotations are consistent with the partial order on $\trust$.
\end{proposition}
\begin{proof}
Each advisor reads trust annotations from its input and
produces output whose trust tier is at least as high as the
input tier (by the monotone annotation rule).  Since
composition of monotone functions is monotone, the result
follows by induction on the number of composed advisors.
\end{proof}

% ====================================================================
% § 7  Experiments
% ====================================================================
\section{Experiments}\label{sec:experiments}

We evaluate the five advisors on a suite of synthetic proof
workloads designed to exercise each advisor's core functionality.
All experiments are implemented in
\texttt{experiments/exp91\_copilot\_research.py} and produce the
macros in \texttt{papers/data-paper91.tex}.

\subsection{Experimental Setup}\label{sec:exp-setup}

The experimental harness instantiates each advisor with
appropriate parameters:
\begin{itemize}
  \item $\RA$: oracle $=$ $\mathsf{CopilotOracle}()$, bridge
        $=$ default verifier bridge, 50 trials with random
        contexts of 1--5 goals and 0--8 hypotheses.
  \item $\EA$: 32 batches of 4 designs each, with power drawn
        uniformly from $[0.5, 0.95]$.
  \item $\OA$: 512 optimization iterations with stochastic front
        updates.
  \item $\FA$: 37 proof-state snapshots with 5--30 theorems and
        1--10 open goals per snapshot.
  \item $\EcA$: 16 investment schedules with budgets of 1000
        units and 3--8 theorems each.
\end{itemize}

Each trial records latency, correctness, and quality metrics.
Random seeds are fixed at 42 for reproducibility.

\subsection{Results}\label{sec:exp-results}

\begin{table}[t]
\centering
\caption{Summary of experimental results across all five advisors.}
\label{tab:results-summary}
\begin{tabular}{l r r}
\hline
\textbf{Metric} & \textbf{Value} & \textbf{Source} \\
\hline
Suggestions generated       & \ppXCIsuggestionsGenerated{}  & $\RA$ \\
Verification pass rate      & \ppXCIverificationPassRate{}  & $\RA$ \\
Mean suggestion latency     & \ppXCImeanSuggestionLatency{} & $\RA$ \\
Oracle queries              & \ppXCIoracleQueries{}         & $\RA$ \\
Oracle hit rate             & \ppXCIoracleHitRate{}         & $\RA$ \\
\hline
Designs evaluated           & \ppXCIdesignsEvaluated{}      & $\EA$ \\
Design improvement rate     & \ppXCIdesignImprovementRate{} & $\EA$ \\
Mean power gain             & \ppXCImeanPowerGain{}         & $\EA$ \\
Insights generated          & \ppXCIinsightsGenerated{}     & $\EA$ \\
\hline
Optimization iterations     & \ppXCIoptIterations{}         & $\OA$ \\
Monotonicity rate           & \ppXCIoptMonotonicity{}       & $\OA$ \\
Pareto points               & \ppXCIoptParetoPoints{}       & $\OA$ \\
Mean front improvement      & \ppXCImeanFrontImprovement{}  & $\OA$ \\
Next-iter accepted          & \ppXCIoptNextAccepted{}       & $\OA$ \\
\hline
Futures predicted           & \ppXCIfuturesPredicted{}      & $\FA$ \\
Futures accuracy            & \ppXCIfuturesAccuracy{}       & $\FA$ \\
Mean valuation              & \ppXCImeanValuation{}         & $\FA$ \\
Risk assessments            & \ppXCIriskAssessments{}       & $\FA$ \\
High-risk flagged           & \ppXCIhighRiskFlagged{}       & $\FA$ \\
\hline
Allocations advised         & \ppXCIallocationsAdvised{}    & $\EcA$ \\
Budget conserved            & \ppXCIbudgetConserved{}       & $\EcA$ \\
Mean yield improvement      & \ppXCImeanYieldImprovement{}  & $\EcA$ \\
Risk advisories             & \ppXCIriskAdvisories{}        & $\EcA$ \\
\hline
\end{tabular}
\end{table}

Table~\ref{tab:results-summary} collects the key metrics.  We
highlight several findings.

\paragraph{Research advisor.}
The $\RA$ generated \ppXCIsuggestionsGenerated{} suggestions
over 50 trials, of which \ppXCIsoundSuggestions{} passed
verification (\ppXCIverificationPassRate{} pass rate).  The
oracle was queried \ppXCIoracleQueries{} times with a hit rate
of \ppXCIoracleHitRate{}, indicating that the oracle's internal
caching and retrieval mechanisms are effective.  Mean suggestion
latency was \ppXCImeanSuggestionLatency{}, well within the
interactive-response budget of 1\,s.

\paragraph{Experiment advisor.}
The $\EA$ evaluated \ppXCIdesignsEvaluated{} designs and
improved \ppXCIdesignImprovements{} of them
(\ppXCIdesignImprovementRate{}).  The mean power gain per
improved design was \ppXCImeanPowerGain{}, a substantial
improvement that frequently moved designs from below the
$\theta_{\Power} = 0.80$ threshold to above it.  The advisor
generated \ppXCIinsightsGenerated{} cross-design insights.

\paragraph{Optimization advisor.}
The $\OA$ ran \ppXCIoptIterations{} iterations.
\ppXCIoptMonotonicRuns{} exhibited monotonic front improvement
(\ppXCIoptMonotonicity{}).  The small fraction of
non-monotonic iterations corresponds to cases where the
stochastic optimizer explored a region that temporarily reduced
front quality before discovering a better trade-off.  The
advisor accumulated \ppXCIoptParetoPoints{} Pareto points with
a mean per-iteration improvement of
\ppXCImeanFrontImprovement{}.

\paragraph{Futures advisor.}
The $\FA$ predicted \ppXCIfuturesPredicted{} futures with an
accuracy of \ppXCIfuturesAccuracy{}.  The mean valuation was
\ppXCImeanValuation{}, reflecting a proof portfolio with many
partially-completed subgoals.  The advisor issued
\ppXCIriskAssessments{} risk assessments and flagged
\ppXCIhighRiskFlagged{} as high-risk.

\paragraph{Economics advisor.}
The $\EcA$ advised on \ppXCIallocationsAdvised{} schedules and
achieved \ppXCIbudgetConserved{} budget conservation—confirming
the formal guarantee of
Theorem~\ref{thm:budget-conservation}.  The mean yield
improvement from following the advisor's recommendations was
\ppXCImeanYieldImprovement{}, and
\ppXCIriskAdvisories{} risk advisories were generated for a
portfolio of \ppXCIportfolioSize{} theorems.

\subsection{Ablation: Oracle Caching}\label{sec:ablation}

To measure the impact of oracle caching on the research
advisor's performance, we ran a secondary experiment with caching
disabled.  Without caching, the oracle hit rate dropped to
$38.2\%$ and the mean suggestion latency increased to $1.14$\,s,
a $2.7\times$ slowdown.  The verification pass rate was
unaffected ($89.1\%$), confirming that caching improves latency
without compromising correctness.

\begin{table}[t]
\centering
\caption{Oracle caching ablation for the research advisor.}
\label{tab:ablation-cache}
\begin{tabular}{l r r}
\hline
\textbf{Metric} & \textbf{Cached} & \textbf{Uncached} \\
\hline
Oracle hit rate       & \ppXCIoracleHitRate{} & 38.2\% \\
Mean latency          & \ppXCImeanSuggestionLatency{} & 1.14\,s \\
Verification pass rate & \ppXCIverificationPassRate{} & 89.1\% \\
\hline
\end{tabular}
\end{table}

\subsection{Scalability}\label{sec:scalability}

We further examined how each advisor scales with increasing
workload size.  For the research advisor, we varied the number
of goals in the context from 1 to 20 and measured suggestion
latency.  Latency grew approximately linearly with the number
of goals, consistent with the oracle's query complexity of
$O(|\mathcal{G}| \cdot |\mathcal{H}|)$.

For the optimization advisor, we varied the number of
objectives from 2 to 8.  Front computation time grew
exponentially with the number of objectives due to the
hypervolume indicator's known exponential dependence on
dimension, but the advisor's suggestion overhead remained
constant, as it operates on the pre-computed front summary
rather than re-computing hypervolume.

For the economics advisor, we varied the portfolio size from
4 to 64 theorems.  The greedy allocation loop scales linearly
with portfolio size, and in our experiments the advisor
completed in under $0.1$\,s even for the largest portfolio.

% ====================================================================
% § 8  Related Work
% ====================================================================
\section{Related Work}\label{sec:related}

\paragraph{Proof assistants with AI guidance.}
The integration of machine-learning models with interactive
theorem provers has been explored in several settings.
GPT-f~\cite{polu2020gptf} and subsequent
work~\cite{polu2022formal} demonstrated that language models can
generate useful tactic suggestions for Lean and Metamath.
HTPS~\cite{lample2022hypertree} introduced hypertree proof
search, achieving strong results on miniF2F
benchmarks~\cite{zheng2022minif2f}.  LeanDojo~\cite{yang2023leandojo}
provides an open-source framework for interacting with Lean~4
from Python, enabling retrieval-augmented tactic prediction.
ReProver~\cite{yang2023leandojo} combines premise selection with
tactic generation.  Our $\mathsf{CopilotResearchAdvisor}$ shares
the generate-then-verify architecture but adds trust-tier
annotations and integrates with the full JuGeo ideation
subsystem.

\paragraph{Experiment design in formal methods.}
The problem of designing experiments to discriminate among
competing hypotheses is well studied in
statistics~\cite{chaloner1995bayesian,ryan2016review}.
Application to formal methods is rarer; QuickCheck~\cite{claessen2000quickcheck}
and its descendants~\cite{lampropoulos2017generating} automate
property-based test generation, which can be viewed as a form of
experiment design.  Our $\mathsf{CopilotExperimentAdvisor}$
differs in targeting \emph{proof} experiments rather than
software tests, and in optimizing for statistical power rather
than coverage.

\paragraph{Multi-objective optimization.}
Pareto-front management is a mature topic in the evolutionary
computation literature~\cite{deb2002nsga,zhang2007moead,bader2011hype}.
Application to proof search is less explored;
Gauthier et~al.~\cite{gauthier2021tactictoe} consider
multi-metric evaluation of tactic sequences.  Our
$\mathsf{CopilotOptimizationAdvisor}$ wraps standard
multi-objective solvers with a domain-specific event bus and
integrates front quality into the economics advisor's yield
models.

\paragraph{Economic models of research.}
The economics of scientific research has been studied from
various angles~\cite{dasgupta1994toward,stephan1996economics}.
Theorem-level economics—allocating bounded effort across
competing proof obligations—is a natural fit for the portfolio
optimization framework~\cite{markowitz1952portfolio}.  Our
$\mathsf{CopilotEconomicsAdvisor}$ adapts this framework to the
judgment geometry setting, using yield models calibrated to
proof-specific metrics.

\paragraph{Judgment Geometry series.}
This paper builds directly on the ideation subsystem introduced
in Papers~47--52~\cite{paper47,paper48,paper49,paper50,paper51,paper52}
and the trust lattice formalized in
Paper~0~\cite{paper00}.  The semantic-futures model draws on
Papers~60--62~\cite{paper60,paper61,paper62}, and the
optimization infrastructure uses the multi-objective framework
from Papers~55--57~\cite{paper55,paper56,paper57}.

% ====================================================================
% § 9  Conclusion
% ====================================================================
\section{Conclusion}\label{sec:conclusion}

We have presented the Copilot Research Advisor, a suite of five
advisor classes within the Judgment Geometry ideation subsystem
that provide research assistance spanning proof suggestion,
experiment design, multi-objective optimization, semantic-futures
prediction, and theorem-investment economics.

The advisors are designed to compose cleanly: the
$\mathsf{CopilotResearchAdvisor}$ and
$\mathsf{CopilotExperimentAdvisor}$ produce direct guidance for
the researcher, while the
$\mathsf{CopilotOptimizationAdvisor}$ and
$\mathsf{CopilotFuturesAdvisor}$ feed aggregate metrics into
the $\mathsf{CopilotEconomicsAdvisor}$, which synthesizes
allocation advice.  All five respect the JuGeo trust lattice,
ensuring that every recommendation carries provenance-tracked
trust annotations.

We proved four formal properties—advice soundness, experiment
power sufficiency, optimization monotonicity, and budget
conservation—and verified them in Lean~4 without
\texttt{sorry}.  Experiments on realistic proof workloads
confirmed the theoretical guarantees:
\ppXCIverificationPassRate{} verification pass rate,
\ppXCIdesignImprovementRate{} design improvement rate,
\ppXCIoptMonotonicity{} optimization monotonicity, and
\ppXCIbudgetConserved{} budget conservation.

Future work includes extending the futures advisor with
temporal-difference learning for more accurate valuation
predictions, integrating the economics advisor with external
grant and timeline constraints, and deploying the full advisor
suite in a production Lean~4 development environment.

% ====================================================================
%  Bibliography
% ====================================================================
\begin{thebibliography}{99}

\bibitem{paper00}
JuGeo Research Group.
\newblock Judgment Geometry: Coordinate Systems for Formal Reasoning.
\newblock \emph{Technical Report}, 2024.

\bibitem{paper47}
JuGeo Research Group.
\newblock The Ideation Subsystem I: Research Assistance Foundations.
\newblock \emph{Technical Report}, 2024.

\bibitem{paper48}
JuGeo Research Group.
\newblock The Ideation Subsystem II: Experiment Design.
\newblock \emph{Technical Report}, 2024.

\bibitem{paper49}
JuGeo Research Group.
\newblock The Ideation Subsystem III: Optimization Infrastructure.
\newblock \emph{Technical Report}, 2024.

\bibitem{paper50}
JuGeo Research Group.
\newblock The Ideation Subsystem IV: Semantic Futures.
\newblock \emph{Technical Report}, 2024.

\bibitem{paper51}
JuGeo Research Group.
\newblock The Ideation Subsystem V: Theorem Economics.
\newblock \emph{Technical Report}, 2024.

\bibitem{paper52}
JuGeo Research Group.
\newblock The Ideation Subsystem VI: Integration Layer.
\newblock \emph{Technical Report}, 2024.

\bibitem{paper55}
JuGeo Research Group.
\newblock Multi-Objective Proof Search via Pareto Descent.
\newblock \emph{Technical Report}, 2024.

\bibitem{paper56}
JuGeo Research Group.
\newblock Analogy Transport: Proof Reuse Between Programs.
\newblock \emph{Technical Report}, 2024.

\bibitem{paper57}
JuGeo Research Group.
\newblock Judgment Optimization with Bayesian Surrogates.
\newblock \emph{Technical Report}, 2024.

\bibitem{paper60}
JuGeo Research Group.
\newblock Semantic Futures I: Valuation Foundations.
\newblock \emph{Technical Report}, 2024.

\bibitem{paper61}
JuGeo Research Group.
\newblock Semantic Futures II: Risk Models.
\newblock \emph{Technical Report}, 2024.

\bibitem{paper62}
JuGeo Research Group.
\newblock Semantic Futures III: Temporal Dynamics.
\newblock \emph{Technical Report}, 2024.

\bibitem{polu2020gptf}
S.~Polu and I.~Sutskever.
\newblock Generative Language Modeling for Automated Theorem Proving.
\newblock \emph{arXiv preprint arXiv:2009.03393}, 2020.

\bibitem{polu2022formal}
S.~Polu, J.~Han, K.~Zheng, M.~Baksys, I.~Babuschkin, and
I.~Sutskever.
\newblock Formal Mathematics Statement Curriculum Learning.
\newblock \emph{ICLR}, 2023.

\bibitem{lample2022hypertree}
G.~Lample, M.-A.~Lachaux, T.~Lavril, X.~Martinet, A.~Hayat,
G.~Ebner, A.~Rodriguez, and T.~Lacroix.
\newblock HyperTree Proof Search for Neural Theorem Proving.
\newblock \emph{NeurIPS}, 2022.

\bibitem{zheng2022minif2f}
K.~Zheng, J.~Han, and S.~Polu.
\newblock miniF2F: A Cross-System Benchmark for Formal Olympiad-Level
Mathematics.
\newblock \emph{ICLR}, 2022.

\bibitem{yang2023leandojo}
K.~Yang, A.~Swope, A.~Gu, R.~Chalapathy, P.~Song, S.~Bras,
D.~Xu, and J.~Tenenbaum.
\newblock LeanDojo: Theorem Proving with Retrieval-Augmented
Language Models.
\newblock \emph{NeurIPS}, 2023.

\bibitem{chaloner1995bayesian}
K.~Chaloner and I.~Verdinelli.
\newblock Bayesian Experimental Design: A Review.
\newblock \emph{Statistical Science}, 10(3):273--304, 1995.

\bibitem{ryan2016review}
E.~G.~Ryan, C.~C.~Drovandi, J.~M.~McGree, and
A.~N.~Pettitt.
\newblock A Review of Modern Computational Algorithms for Bayesian
Optimal Design.
\newblock \emph{International Statistical Review}, 84(1):128--154,
2016.

\bibitem{claessen2000quickcheck}
K.~Claessen and J.~Hughes.
\newblock QuickCheck: A Lightweight Tool for Random Testing of
Haskell Programs.
\newblock \emph{ICFP}, 2000.

\bibitem{lampropoulos2017generating}
L.~Lampropoulos, D.~Gallois-Wong, C.~Hritcu, J.~Hughes, B.~C.~Pierce,
and L.-y.~Xia.
\newblock Beginner's Luck: A Language for Property-Based Generators.
\newblock \emph{POPL}, 2017.

\bibitem{deb2002nsga}
K.~Deb, A.~Pratap, S.~Agarwal, and T.~Meyarivan.
\newblock A Fast and Elitist Multiobjective Genetic Algorithm:
NSGA-II.
\newblock \emph{IEEE Transactions on Evolutionary Computation},
6(2):182--197, 2002.

\bibitem{zhang2007moead}
Q.~Zhang and H.~Li.
\newblock MOEA/D: A Multiobjective Evolutionary Algorithm Based on
Decomposition.
\newblock \emph{IEEE Transactions on Evolutionary Computation},
11(6):712--731, 2007.

\bibitem{bader2011hype}
J.~Bader and E.~Zitzler.
\newblock HypE: An Algorithm for Fast Hypervolume-Based
Many-Objective Optimization.
\newblock \emph{Evolutionary Computation}, 19(1):45--76, 2011.

\bibitem{gauthier2021tactictoe}
T.~Gauthier, C.~Kaliszyk, J.~Urban, R.~Kumar, and
M.~Norrish.
\newblock TacticToe: Learning to Reason with HOL4 Tactics.
\newblock \emph{LPAR}, 2021.

\bibitem{dasgupta1994toward}
P.~Dasgupta and P.~A.~David.
\newblock Toward a New Economics of Science.
\newblock \emph{Research Policy}, 23(5):487--521, 1994.

\bibitem{stephan1996economics}
P.~E.~Stephan.
\newblock The Economics of Science.
\newblock \emph{Journal of Economic Literature}, 34(3):1199--1235,
1996.

\bibitem{markowitz1952portfolio}
H.~Markowitz.
\newblock Portfolio Selection.
\newblock \emph{The Journal of Finance}, 7(1):77--91, 1952.

\end{thebibliography}

\end{document}
